\section{工程结论、局限与课堂问答}

最后一组幻灯片把论文结果压缩成工程原则：奖励中间 thought、使用 dense signal、检查 token efficiency，并把 curriculum、front-loading 与 reinforcement 组合起来。但真正可落地的结论仍需经过数据审计、cost accounting、failure analysis 与安全边界。

\subsection{三项 ablation：为什么有效、信号多密、数据多省}

Key Ablations 页把 RLP 的价值拆成三点：Reward to Think 检验 thought 是否获得直接 credit；Sparse vs Dense 检验只奖励最终正确与逐位置信息增益的差别；Token Efficiency 检验相对超大 NTP baseline 的数据利用率。三列共同说明收益来自 objective design，而不只是多生成几段 CoT。

\lecturefigure{slide-041.jpg}{RLP 的 reward-to-think、dense reward 与 token-efficiency 结论}{Stanford 官方录像，00:45:56--00:46:27}

\begin{warningbox}{Ablation 仍需完整 cost table}
若 dense reward 使用更多 rollout、更多 scorer forward 或更长 thought，它同时改变 signal 与 compute。理想 ablation 应报告相同 FLOPs、相同 wall-clock、相同 token 或多种匹配方式，避免把额外计算误当作奖励形式本身的效果。
\end{warningbox}

\subsection{Pascal 到 Hopper：能力来自策略叠加}

曲线把 four learners 重新连起来：curriculum 带来第一层改善，front-loading reasoning 再提高 foundation，RLP 加入 learning through thinking 后达到最高相对提升。曲线是跨实验的教学综合，不是一次统一训练中严格可加的 causal decomposition，因此不能把每段百分比当成独立可叠加常数。

\lecturefigure{slide-042.jpg}{从 Pascal 到 Hopper：逐步加入三种学习策略}{Stanford 官方录像，00:46:27--00:47:12}

\begin{importantbox}{综合判断而不是机械相加}
Curriculum 改变模型看到的分布，front-loading 改变 base representation，RLP 又改变 thought policy 的优化地形；三者存在交互。真正组合时应做 factorial ablation，而不是把各论文 headline gain 相加预测最终收益。
\end{importantbox}

\subsection{六条结论与三个非结论}

课程最终总结：two-phase 有效；front-loading 产生 durable and compounding advantage；高质量信号可能在 SFT 后显现；RLP 把 RL reasoning 前移；普通未标注文本也能提供 reasoning-like training signal；探索与奖励为 reasoning scaling 打开新轴。每条都应附带模型、数据和评测边界。

\lecturefigure{slide-043.jpg}{Curriculum、reasoning 与 reinforcement 重定义预训练}{Stanford 官方录像，00:47:12--00:49:00}

图中\textbf{没有证明}三件事：第一，所有 chain-of-thought 都值得保留；第二，RLP 会自动减少 hallucination 或提高 alignment；第三，language-only 结果会直接迁移到 vision-language model。讲者在 Q\&A 中明确把 multimodal 视为未测试，把 RLHF/hallucination 视为相邻但不属于本工作的议题。

\subsection{最终 recipe 回到系统协作}

闭幕页再次出现四组件，但此时每一列已有具体含义：Smart Data 需要可追踪的 quality/domain labels；Smart Architecture 需要支持长训练与 rollout；Smart Algorithms 包含 curriculum、front-loading 与 RLP；Smart Collaboration 负责统一 pretraining/post-training 指标和 cost budget。它提醒读者：算法论文只有进入完整 pipeline 才能转化为模型能力。

\lecturefigure{slide-044.jpg}{完成机制学习后重新审视四组件 recipe}{Stanford 官方录像，00:49:00--00:49:30}

\begin{knowledgebox}{一个可执行的决策顺序}
先审计数据与评测口径，再用小规模 proxy 估计 quality / repeat / phase boundary；随后在固定预算下测试 front-loading；只有当 baseline、FLOPs 与后训练流程稳定后，再引入 RLP rollout。每一步都保留未采用方案和失败结果，避免把 recipe 调参变成不可复核的炼金术。
\end{knowledgebox}

\subsection{课堂问答：实现细节与边界}

问答补充了多个关键点。RLP 的 advantage machinery 与 GRPO 相似，主要变化是 information-gain reward；实验在文档中随机选 token，而不是基于 entropy 精筛；RLP 使用一般 pretraining mixture，包括大量 web crawl，并非只在高质量 reasoning data 上运行；从约完整预训练 20\% 的更早 checkpoint 开始也观察到相对 CPT 的收益。

\teachervoice{00:49:54--00:50:35，讲者说明 RLP 复用了 GRPO 式 group-relative advantage，创新重点是 reward construction。}

\teachervoice{00:55:36--00:57:43，讲者纠正提问者：RLP 不是只在最高质量数据上做；实践中随机选择 token，并且在更早的约 4T-token checkpoint 上也做过有利实验。}

\begin{lstlisting}[caption={评估 RLP 时建议记录的最小实验表}]
experiment = {
    "start_checkpoint_tokens": ...,   # e.g. 19.8T
    "extra_data_tokens": ...,
    "rollouts_per_position": ...,
    "mean_thought_length": ...,
    "estimated_training_flops": ...,
    "gpu_hours": ...,
    "token_selection": "uniform_random",
    "data_mixture_version": ...,
    "post_training_recipe": ...,
    "benchmark_harness": ...,
}
\end{lstlisting}

\subsection{数据质量、RLHF 与 multimodal 的边界}

对于一般网页质量，讲者提到 FineWeb-Edu 式 educational classifier：用 rubric 给文档 1--5 分，并可扩展到 domain 与教育层级。对于 RLHF bias、hallucination 与 reward hacking，她明确表示这是 alignment 的持续研究，不是本讲三项工作的直接结论。对于 vision-language alignment，她的回答是尚未测试。

\begin{warningbox}{不要把问答中的谨慎回答改写成确定结论}
“RLHF 仍可能占 post-training 的一小部分”“classifier 需要持续更新”“RLAIF 可能降低部分主观性”都是课堂讨论与经验判断，不是本讲实验验证。讲义保留其边界，避免把相邻领域观点写成 RLP 的效果声明。
\end{warningbox}

\subsection{本章小结}

工程上最重要的不是复述 headline gain，而是建立完整控制面：data quality 与 curriculum、reasoning allocation、rollout reward、EMA stability、token/FLOP/wall-clock matching、post-training persistence 和跨域限制。只有这些信息齐全，才能判断 RLP 或 front-loading 是否值得进入生产训练。

